\documentclass[11pt,a4paper]{article}
\usepackage[italian]{babel}
\usepackage[margin=2.5cm]{geometry}
\usepackage{parskip}
\usepackage{amsmath, amssymb, amsthm}
\usepackage{booktabs}
\usepackage{array}
\usepackage{xcolor}
\usepackage{needspace}
\usepackage{enumitem}
\usepackage{listings}
\usepackage{tikz}
\usepackage[most]{tcolorbox}
\usepackage[hidelinks]{hyperref}
\usetikzlibrary{decorations.pathreplacing, shapes.geometric, arrows.meta, positioning}

% Niente vedove/orfane: i paragrafi non lasciano righe isolate a fine/inizio pagina
\widowpenalty=10000
\clubpenalty=10000

% Elenchi più compatti
\setlist{itemsep=2pt, parsep=0pt, topsep=4pt, partopsep=0pt}

% --- Riquadri con bordo nero, senza numero (si spezzano tra due pagine) ---
% Uso: \begin{definizione} ... \end{definizione}
%      \begin{teorema}[Nome] ... \end{teorema}  (il nome è facoltativo)
\tcbset{nota/.style={enhanced jigsaw, breakable, colback=white,
  colframe=black, boxrule=0.5pt, sharp corners}}
\NewTColorBox{definizione}{}{nota,
  before upper={\textbf{Definizione.}\ }}
\NewTColorBox{teorema}{o}{nota,
  before upper={\textbf{Teorema\IfValueT{#1}{ (#1)}.}\ }}
\NewTColorBox{proposizione}{o}{nota,
  before upper={\textbf{Proposizione\IfValueT{#1}{ (#1)}.}\ }}
\NewTColorBox{proprieta}{o}{nota,
  before upper={\textbf{Proprietà\IfValueT{#1}{ (#1)}.}\ }}
\NewTColorBox{assioma}{o}{nota, boxrule=1pt,
  before upper={\textbf{Assioma\IfValueT{#1}{ (#1)}.}\ }}

% --- Ambienti semplici (senza riquadro) ---
% Il titolo sta in cima al blocco, anche se il contenuto inizia con un riquadro o
% con codice; e non resta mai da solo in fondo a una pagina.
% Uso: \begin{esempio}[Titolo facoltativo] ... \end{esempio}
\newcommand{\newrunin}[2]{%
  \NewDocumentEnvironment{#1}{o}{%
    \par\Needspace{4\baselineskip}\addvspace{\medskipamount}\noindent
    \textbf{#2}\IfValueT{##1}{ (##1)}\textbf{.}\ \ignorespaces}%
  {\par\addvspace{\medskipamount}}}
\newrunin{esempio}{Esempio}
\newrunin{esercizio}{Esercizio}
\newrunin{osservazione}{Osservazione}

% V e F nelle tavole di verità
\newcommand{\V}{\textsf{V}}
\newcommand{\F}{\textsf{F}}

% --- Codice C ---
% Uso: \begin{codice} ... \end{codice}      codice C numerato
%      \begin{comando} ... \end{comando}    comandi da terminale
%      \begin{risultato} ... \end{risultato} output di un programma
%      \lstinline|printf("ciao");|           codice dentro al testo
\lstdefinestyle{c}{
  language=C,
  basicstyle=\ttfamily\small,
  keywordstyle=\color{blue}\bfseries,
  commentstyle=\color{gray}\itshape,
  stringstyle=\color{red!70!black},
  numbers=left,
  numberstyle=\tiny\color{gray},
  numbersep=8pt,
  columns=flexible,
  keepspaces=true,
  breaklines=true,
  showstringspaces=false,
  tabsize=4,
  morekeywords={include, define},
  % lettere accentate nei commenti
  literate={à}{{\`a}}1 {è}{{\`e}}1 {é}{{\'e}}1 {ì}{{\`i}}1
           {ò}{{\`o}}1 {ù}{{\`u}}1 {È}{{\`E}}1 {À}{{\`A}}1
}
\lstset{style=c}
\newtcblisting{codice}{listing only, nota, left=6mm,
  listing options={style=c}}
\newtcblisting{comando}{listing only, nota,
  listing options={basicstyle=\ttfamily\small, numbers=none, language={},
    columns=flexible, keepspaces=true, breaklines=true}}
\newtcblisting{risultato}{listing only, enhanced jigsaw, breakable,
  colback=black!5, colframe=black, boxrule=0.5pt, sharp corners,
  listing options={basicstyle=\ttfamily\small, numbers=none, language={},
    columns=flexible, keepspaces=true, breaklines=true}}

% --- Diagrammi di flusso (simboli standard) ---
\tikzset{
  terminale/.style = {ellipse, draw, minimum width=2.5cm, minimum height=0.9cm},
  io/.style        = {trapezium, trapezium left angle=70, trapezium right angle=110,
                      draw, minimum width=2.5cm, minimum height=0.9cm},
  processo/.style  = {rectangle, draw, minimum width=2.5cm, minimum height=0.9cm},
  decisione/.style = {diamond, draw, aspect=2, minimum width=2.5cm},
  freccia/.style   = {thick, -{Stealth}}
}

\title{Funzioni suriettive, iniettive e biunivoche}
\author{Marco Cerbella}
\date{Lezione 5 -- 8 ottobre 2026}

\begin{document}
\maketitle

\section{Funzioni suriettive}

\begin{definizione}
Una applicazione $f : D \subseteq \mathbb{R} \to C \subseteq \mathbb{R}$ si dice \emph{suriettiva} se
\[
\forall y \in C, \ \exists x \in D \text{ tale che } f(x) = y.
\]
\end{definizione}

Significato geometrico: ogni retta parallela all'asse delle ascisse che taglia l'asse delle ordinate in un punto $y \in C$ interseca il grafico di $f$ in \emph{almeno} un punto.

\begin{esempio}
\leavevmode
\begin{itemize}
    \item $f : \mathbb{R} \to \mathbb{R}$, $f(x) = ax + b$ (con $a \neq 0$) è suriettiva perché $f(\mathbb{R}) = \mathbb{R}$.
    \item $g : \mathbb{R} \to \mathbb{R}$, $g(x) = x^2$ \emph{non} è suriettiva perché $g(\mathbb{R}) = [0, +\infty) \subsetneq \mathbb{R}$: un numero negativo non può essere il quadrato di un numero reale.
\end{itemize}
\end{esempio}

\section{Funzioni iniettive}

\begin{definizione}
Una applicazione $f : D \to \mathbb{R}$ si dice \emph{iniettiva} se
\[
\forall x_1, x_2 \in D, \quad x_1 \neq x_2 \Rightarrow f(x_1) \neq f(x_2).
\]
\end{definizione}

Significato geometrico: ogni retta parallela all'asse delle ascisse che taglia l'asse delle ordinate in un punto $y \in C$ interseca il grafico di $f$ in \emph{al più} un punto.

\begin{esempio}
\leavevmode
\begin{itemize}
    \item $f(x) = ax + b$ è iniettiva per $a \neq 0$. Infatti, per assurdo, siano $x_1 \neq x_2$ con $f(x_1) = f(x_2)$, cioè $a x_1 + b = a x_2 + b$. Sottraendo $b$ e dividendo per $a$ si ottiene $x_1 = x_2$, contro l'ipotesi.
    \item $g(x) = x^2$ \emph{non} è iniettiva: per $x_1 \neq 0$ si ha $x_1 \neq -x_1$ ma $g(x_1) = x_1^2 = g(-x_1)$.
\end{itemize}
\end{esempio}

\section{Funzioni biunivoche}

\begin{definizione}
$f$ si dice \emph{biunivoca} se è sia iniettiva sia suriettiva.
\end{definizione}

Significato geometrico: ogni retta parallela all'asse delle ascisse che taglia l'asse delle ordinate in un punto $y \in C$ interseca il grafico di $f$ in \emph{uno ed un solo} punto.

\begin{esempio}
$f(x) = ax + b$ è biunivoca per $a \neq 0$; $g(x) = x^2$ (da $\mathbb{R}$ in $\mathbb{R}$) non è biunivoca.
\end{esempio}

\end{document}
